\chapter{Two local directions}
\label{ch:two-directions}

\section{The phrase that named two objects}

A protected field theory on $X^{\mathrm{an}}\times\mathbb R$, with $X$ a
smooth complex algebraic curve, has two elementary local motions. Points
collide in the holomorphic curve and produce a singular operator product.
Intervals compose along the oriented real line and produce an ordered
associative product. Each motion has its own associative structure, and
the phrase \emph{noncommutative chiral algebra} has been used for both.

\begin{definition}[Transversely ordered chiral algebra]
\label{def:transverse}
A transversely ordered chiral algebra is an augmented $E_1$-algebra for
the componentwise coefficient tensor in factorisable right $D$-modules,
\[
  A \in \Alg^{\mathrm{aug}}_{E_1}\bigl(\FactD(X), \otimes^!\bigr).
\]
The coefficient records collision on $X$; the internal $E_1$
multiplication records ordered fusion along the transverse line. The
product $\otimes^!$ preserves a fixed finite support.
\end{definition}

\begin{definition}[Associative chiral algebra]
\label{def:associative-chiral}
An associative chiral algebra is an augmented $E_1$-algebra for chiral
convolution,
\[
  C \in \Alg^{\mathrm{aug}}_{E_1}\bigl(\DMod(\Ran X), \starch\bigr),
\]
equivalently an algebra over the non-symmetric associative chiral operad.
The product $\starch$ assembles disjoint supports.
\end{definition}

These are different structures over different monoidal products, and
neither determines the other. Definition~\ref{def:transverse} is a theory
on $X\times\mathbb R$; Definition~\ref{def:associative-chiral} is a
theory on $X$ alone with the symmetry of the chiral product dropped.

\begin{proposition}[Formal-disc realization]
\label{prop:nonlocal-va}
\StatusConditional{} On the formal disc, the field realization of an
associative chiral algebra is a nonlocal vertex algebra in the sense of
Li \cite{Li05}, equivalently a field algebra in the sense of
Bakalov--Kac \cite{BK03}: the state-field correspondence and weak
associativity are retained, and locality is not imposed. An exchange
operator $S(z)$ satisfying unitarity, the spectral Yang--Baxter equation,
and braided locality upgrades the realization to a quantum vertex algebra
in the sense of Etingof--Kazhdan \cite{EK00}.
\end{proposition}

The hypotheses of Proposition~\ref{prop:nonlocal-va} are the
identification of the formal-disc coefficient category with the category
of fields, and the compatibility of $\starch$ with the state-field
correspondence. The braided upgrade is additional data: an associative
chiral algebra carries no canonical $S(z)$.

\section{Three bars, not one}

Each structure has its own bar construction, and the three constructions
below contract different geometric events.

\begin{notation}[The three bars]
\label{not:three-bars}
$\Barperp$ is the internal associative bar for the transverse product
$\otimes^!$; its differential multiplies adjacent transverse letters.
$\BarchAss$ is the operadic bar of the associative chiral product; its
differential contracts chiral composition trees and retains the whole
diagonal distribution. $\BarchLie$ is chiral Chevalley chains of a chiral
Lie algebra, the object of Francis--Gaitsgory chiral Koszul duality
\cite{FG12}; it is neither of the preceding two.
\end{notation}

The underlying gradings of $\Barperp$ and $\BarchAss$ count different
geometric events: subdivisions of an interval in the first case,
collision trees on the curve in the second. An asserted equality between
them is meaningless without a functor carrying one multiplication problem
to the other; see No-go~\ref{nogo:unnamed-bar-comparison}.

\section{The doubly noncommutative object}

\begin{definition}[Doubly associative chiral algebra]
\label{def:doubly}
A doubly associative chiral algebra is a triple $(A, m_\perp, m_X,
\lambda)$ in which $m_\perp$ is an augmented transverse $E_1$
multiplication for $\otimes^!$, $m_X$ is an augmented associative chiral
$E_1$ multiplication for $\starch$, and $\lambda$ is a mixed distributive
law between them, including all higher homotopies. Completions in the
nilpotent, weight, pole, coradical, and analytic filtrations are
specified separately.
\end{definition}

Any subset of this datum defines the corresponding partial theory. No
missing component is inferred from the others; in particular $\lambda$ is
supplied, never deduced from $m_\perp$ and $m_X$.

\begin{proposition}[Where the mixed law can live]
\label{prop:lambda-sector}
\StatusProved{} The sector on which a mixed distributive law between
$\otimes^!$ and $\starch$ is geometrically defined is the
common-decomposition sector, that is the image of the idempotent $e$ of
Theorem~\ref{thm:aligned-splitting}. Off that sector the two products
admit no interchange, only the retraction of
Theorem~\ref{thm:aligned-splitting}, so no distributive law can be
written.
\end{proposition}

\begin{proof}
A distributive law requires a natural transformation interchanging the
two products in a fixed direction, subject to the coherence conditions of
each. By Theorem~\ref{thm:aligned-splitting} the canonical comparison in
the additive model is $\iota$ with retraction $\pi$, and $\pi\iota=\id$
while $\iota\pi=e\ne\id$ in general. A natural transformation subject to
both coherence conditions therefore exists only where $e$ acts as the
identity, which is $\operatorname{im}e$ by
Corollary~\ref{cor:aligned-duoidal}.
\end{proof}

\begin{remark}[Open: an instance of $\lambda$]
\label{rem:lambda-open}
Proposition~\ref{prop:lambda-sector} locates $\lambda$ but does not
construct one. No example in the present corpus exhibits a mixed
distributive law explicitly. Zamolodchikov--Faddeev algebras are the
natural candidate, since their spectral variable belongs to the chiral
direction while their word belongs to the transverse direction; producing
$\lambda$ for a Zamolodchikov--Faddeev algebra, and verifying that it
lands in $\operatorname{im}e$, is the outstanding construction on which
Definition~\ref{def:doubly} depends for content.
\end{remark}

\section{The order of dependence}

The remaining structures of the theory are functorial extensions with
distinct sources, and none is determined by its predecessor. A second
transverse direction, or a flat meromorphic connection on
two-dimensional configurations, supplies braid transport. A cyclic trace
supplies the circle object and ribbon contractions. A
clutching-compatible modular action supplies sewing over stable curves.
The transverse ribbon genus $g_\perp$ and the Deligne--Mumford chiral
genus $g_X$ are independent variables.

Scalars occur at specified points of
\[
  \FactD(X) \longrightarrow A \longrightarrow \Barperp(A)
  \longrightarrow Z(A) \longrightarrow \mathrm{Rep}(H)
  \longrightarrow K_0 \longrightarrow \text{character},
\]
and an equality of scalars is explained only by a comparison morphism
between the structured objects preceding them.
